\clearpage
\section*{Заключение}
\addcontentsline{toc}{section}{Заключение}

В ходе первой недели учебной практики последовательно выполнены пять заданий, охватывающих методы расширения, LINQ, итераторы, интерфейсы и рефлексию. Для каждого задания создана отдельная ветка репозитория, библиотека классов и проект тестирования. Решение собирается автоматически, а корректность подтверждается 32 модульными тестами.

Рефлексивный анализ показал, что наиболее существенные затруднения возникали не при записи отдельных операторов языка, а при уточнении контрактов и границ поведения. К таким вопросам относятся обработка пустых данных, выбор уровня агрегации, понимание отложенного выполнения, наблюдаемость операций с возвращаемым типом \texttt{void} и настройка области поиска средствами рефлексии. Во всех случаях решение было получено через декомпозицию задачи, формализацию ожидаемого результата и проверку отдельными тестами.

Практическая работа подтвердила, что модульные тесты выполняют одновременно контрольную, регрессионную и документирующую функции. Они увеличивают объём начальной работы, однако сокращают время поиска ошибок и обеспечивают безопасное развитие решения. Автоматический запуск тестов в CI дополняет локальную проверку и подтверждает воспроизводимость сборки.

Поставленная цель достигнута: результаты работы систематизированы, трудности и способы их преодоления проанализированы, полный набор модульных тестов описан через проверяемые случаи, предусловия и постусловия. Полученные навыки создают основу для дальнейшего изучения архитектуры приложений, метаданных, динамических библиотек и автоматизированного контроля качества программного обеспечения.
